\chapter{Размерностная трансмутация и планковская самосогласованность}
\label{ch:v10-induced-newton-dimensional-transmutation}

\section{Две различные бета-величины}

Необходимо различать однопетлевой коэффициент $b$ функции
$\beta_g=\mu\,dg/d\mu$ и безразмерный коэффициент $\beta_E$ в
$B=\beta_E m$. Совпадение названий не является физической связью.

Ранний relative-$U(1)$ сектор строго дал
\begin{equation}
 b=2,
 \qquad g^2(\mu_{\rm spec})=\frac38.
 \label{eq:v10-newton-dt-inherited-data}
\end{equation}
Поэтому однопетлевое отношение Ландау равно
\begin{equation}
 \log\frac{\Lambda_L}{\mu_{\rm spec}}
 =\frac{8\pi^2}{b g^2(\mu_{\rm spec})}
 =\frac{32\pi^2}{3}.
 \label{eq:v10-newton-dt-landau-ratio}
\end{equation}
Для семени обратной площади $m_{\rm DT}=\Lambda_L^2$ получается
\begin{equation}
 \log\frac{m_{\rm DT}}{\mu_{\rm spec}^2}
 =\frac{64\pi^2}{3}.
 \label{eq:v10-newton-dt-seed-ratio}
\end{equation}
Это точное безразмерное отношение, но при положительном $b$ оно описывает
ультрафиолетовый полюс Ландау, а не инфракрасную гравитационную щель.
Кроме того, перенос этих данных в текущий ньютоновский носитель не выведен.

\section{Планковская самосогласованность}

Ранняя метафора локальных планковских часов соответствует условию
\begin{equation}
 m g_N=1,
 \qquad g_N=\frac{\hbar G}{c^3}.
 \label{eq:v10-newton-dt-planck-condition}
\end{equation}
Вместе с эйнштейновским равенством
\begin{equation}
 16\pi\beta_E m g_N=1
 \label{eq:v10-newton-dt-einstein-condition}
\end{equation}
оно выбирает только
\begin{equation}
 \boxed{\beta_E=\frac1{16\pi}}.
 \label{eq:v10-newton-dt-beta-E}
\end{equation}
Размерная величина сокращается. Следовательно, планковская самоссылка может
условно выбрать безразмерную нормировку, но не абсолютное значение $G$.

\section{Условный общий родитель}

Введём нормированные переменные
\begin{equation}
 u_{\rm RG}=\frac{m}{\mu_{\rm spec}^2\exp(64\pi^2/3)},
 \qquad u_P=m g_N,
 \qquad u_E=16\pi\beta_E m g_N.
\end{equation}
Они собираются в положительный родитель
\begin{equation}
 \mathcal P_{\rm DT}
 =\frac12(u_{\rm RG}-1)^2
 +\frac12(u_P-u_{\rm RG})^2
 +\frac12(u_E-u_P)^2.
 \label{eq:v10-newton-dt-parent}
\end{equation}
Его единственный нуль равен $u_{\rm RG}=u_P=u_E=1$. Гессиан
\begin{equation}
 \begin{pmatrix}2&-1&0\\-1&2&-1\\0&-1&1\end{pmatrix}
\end{equation}
имеет ранг $3$, определитель $1$ и ведущие миноры $(2,3,1)$.

Однако в логарифмических переменных $(m,\mu_{\rm spec}^2,g_N)$ размерные
условия имеют матрицу
\begin{equation}
 C_{\rm DT}=
 \begin{pmatrix}1&-1&0\\1&0&1\end{pmatrix},
 \qquad \rank C_{\rm DT}=2,
 \qquad C_{\rm DT}(1,1,-1)^T=0.
 \label{eq:v10-newton-dt-scale-map}
\end{equation}
Одна общая масштабная мода сохраняется. Независимая фиксация $g_N$ повысила
бы ранг до $3$, но являлась бы внешним ньютоновским входом.

\begin{theorem}[Планковская самосогласованность выбирает коэффициент]
При условии $m g_N=1$ эйнштейновская нормировка выбирает
$\beta_E=1/(16\pi)$. Этот результат безразмерен и не зависит от значения
$G$.
\end{theorem}

\begin{nogocriterion}[RG-отношение не является ньютоновским масштабом]
Положительная однопетлевая функция с $b=2$ даёт ультрафиолетовое отношение
Ландау, а не инфракрасную массу. Даже после условного планковского замыкания
матрица размерных условий сохраняет ядро $(1,1,-1)$. Поэтому абсолютный
$G$ не выведен.
\end{nogocriterion}

\begin{protocol}[Статус трансмутации и самосогласованности]\begin{enumerate}
\item исторические RG-данные: \textbf{$2/2$};
\item перенос RG-данных в текущий носитель: \textbf{$0/2$};
\item условная архитектура: \textbf{$10/10$};
\item условное замыкание: \textbf{$7/7$};
\item выбор $\beta_E=1/(16\pi)$: \textbf{условно $1/1$};
\item физическое происхождение входов: \textbf{$0/4$};
\item абсолютные $m$ и $G$: \textbf{$0/2$};
\item следующий гейт: \textbf{аномальное дыхание сквозного потока}.
\end{enumerate}\end{protocol}

Машинный сертификат сохранён в
\path{s2t_v10_cell_birth_four_volume_induced_newton_dimensional_transmutation_beta_parent_origin_gate_results.json}.